\documentclass[preview]{standalone}

\usepackage{amsmath}
\usepackage{amssymb}
\usepackage{stellar}
\usepackage{definitions}

\begin{document}

\id{quadratic-integers-ideal-factorization}
\genpage

\section{Prime ideals}

\begin{snippetdefinition}{sqrt-minus-five-special-ideals-definition}{Four ideals in the integers with square root of minus five}
    Let \(A=\integers[\sqrt{-5}]\). Define
    \[
        I_1=(2,1-\sqrt{-5}),
        \qquad
        I_2=(2,1+\sqrt{-5}),
    \]
    and
    \[
        J_1=(3,1-\sqrt{-5}),
        \qquad
        J_2=(3,1+\sqrt{-5}).
    \]
\end{snippetdefinition}

\begin{snippetproposition}{sqrt-minus-five-special-ideals-maximal}{The four ideals are maximal}
    The \ideal[ideals]
    \(I_1,I_2,J_1,J_2\) of
    \(\integers[\sqrt{-5}]\) are maximal, and hence prime.
\end{snippetproposition}

\begin{snippetproof}{sqrt-minus-five-special-ideals-maximal-proof}{sqrt-minus-five-special-ideals-maximal}{The four ideals are maximal}
    Put \(A=\integers[\sqrt{-5}]\). Define the surjective
    \ringhomomorphism
    \[
        \varphi\colon A\fromto\integers/2\integers,
        \qquad
        \varphi(a+b\sqrt{-5})=[a+b]_2.
    \]
    It respects multiplication because, modulo \(2\), substituting
    \(\sqrt{-5}=1\) is compatible with
    \((\sqrt{-5})^2=-5\equiv1\pmod2\). Moreover,
    \[
        \varphi(a+b\sqrt{-5})=0
        \ifandonlyif
        a+b\equiv0\pmod2.
    \]
    Since
    \[
        a+b\sqrt{-5}
        =
        (a+b)+b(\sqrt{-5}-1),
    \]
    this condition is equivalent to membership in \(I_1\). Thus
    \[
        A/I_1\ringisomorphic\integers/2\integers.
    \]
    The quotient is a \field, so \(I_1\) is maximal.

    For \(I_2\), use
    \[
        a+b\sqrt{-5}\longmapsto[a-b]_2.
    \]
    In fact \(I_1=I_2\), because their second generators differ by
    \(2\sqrt{-5}\in(2)\).

    Similarly, the \ringhomomorphism
    \[
        a+b\sqrt{-5}\longmapsto[a+b]_3
    \]
    has kernel \(J_1\), whereas
    \[
        a+b\sqrt{-5}\longmapsto[a-b]_3
    \]
    has kernel \(J_2\). Hence
    \[
        A/J_1\ringisomorphic\integers/3\integers,
        \qquad
        A/J_2\ringisomorphic\integers/3\integers.
    \]
    Both quotients are \field[fields], so \(J_1\) and \(J_2\) are
    maximal.
\end{snippetproof}

\section{Factorization of principal ideals}

\begin{snippetproposition}{sqrt-minus-five-principal-ideal-factorizations}{Factorizations of the ideals generated by two, three, and six}
    In \(A=\integers[\sqrt{-5}]\),
    \[
        (2)=I_1I_2,
        \qquad
        (3)=J_1J_2,
    \]
    and consequently
    \[
        (6)=I_1I_2J_1J_2.
    \]
    Since \(I_1=I_2\), the first equality may also be written
    \((2)=I_1^2\).
\end{snippetproposition}

\begin{snippetproof}{sqrt-minus-five-principal-ideal-factorizations-proof}{sqrt-minus-five-principal-ideal-factorizations}{Factorizations of the ideals generated by two, three, and six}
    Products of finitely generated \ideal[ideals] are generated by the
    pairwise products of their generators. Hence
    \begin{align*}
        I_1I_2
        &=
        (4,\ 2(1+\sqrt{-5}),\ 2(1-\sqrt{-5}),\\
        &\hspace{4.2cm}
        (1-\sqrt{-5})(1+\sqrt{-5}))\\
        &=
        (4,\ 2(1+\sqrt{-5}),\ 2(1-\sqrt{-5}),\ 6).
    \end{align*}
    Every generator belongs to \((2)\), so \(I_1I_2\subseteq(2)\).
    Conversely, \(2=6-4\in I_1I_2\), proving \((2)=I_1I_2\).

    Similarly,
    \[
        J_1J_2
        =
        (9,\ 3(1+\sqrt{-5}),\ 3(1-\sqrt{-5}),\ 6).
    \]
    This ideal is contained in \((3)\), while \(3=9-6\) belongs to the
    displayed generating set. Therefore \((3)=J_1J_2\). Finally,
    \[
        (6)=(2)(3)=I_1I_2J_1J_2.
    \]
\end{snippetproof}

\begin{snippet}{sqrt-minus-five-element-ideal-factorization-contrast}
    The principal \ideal \((6)\) factors as a product of prime ideals,
    even though the element \(6\) does not factor as a product of prime
    elements in \(\integers[\sqrt{-5}]\).
\end{snippet}

\section{Dimension-one behavior}

\begin{snippetproposition}{sqrt-minus-five-nonzero-prime-ideals-maximal}{Nonzero prime ideals are maximal}
    Every nonzero \primeideal of
    \[
        A=\integers[\sqrt{-5}]
    \]
    is a \maximalideal.
\end{snippetproposition}

\begin{snippetproof}{sqrt-minus-five-nonzero-prime-ideals-maximal-proof}{sqrt-minus-five-nonzero-prime-ideals-maximal}{Nonzero prime ideals are maximal}
    Let \(I\idealin A\) be a nonzero \primeideal. Choose
    \(\alpha=a+b\sqrt{-5}\in I\difference\{0\}\). Then
    \[
        N(\alpha)=\alpha\overline{\alpha}
        \in I\intersection\integers
    \]
    is nonzero. The contraction \(I\intersection\integers\) is a proper
    prime \ideal of \(\integers\), and hence
    \[
        I\intersection\integers=(p)
    \]
    for a \primen \(p\).

    Since \(I\) is prime, \(A/I\) is an \integraldomain. Let
    \[
        \alpha+I\neq0_{A/I},
        \qquad
        \alpha=a+b\sqrt{-5},
    \]
    and put \(\beta=\overline{\alpha}\).

    If \(p\nmid N(\alpha)\), Bézout's identity gives \(x,y\in\integers\)
    with
    \[
        xN(\alpha)+yp=1.
    \]
    Since \(p\in I\),
    \[
        (\alpha+I)(x\beta+I)
        =
        xN(\alpha)+I
        =
        1+I.
    \]

    Suppose instead that \(p\divides N(\alpha)\). Then
    \(\alpha\beta=N(\alpha)\in I\). Since \(\alpha\notin I\) and \(I\)
    is prime, \(\beta\in I\). Therefore
    \[
        \alpha
        =
        \beta+2b\sqrt{-5}
        \equiv
        2b\sqrt{-5}\pmod I.
    \]
    The class is nonzero, so \(p\nmid2b\); otherwise
    \(2b\sqrt{-5}\in(p)\subseteq I\). Also \(p\nmid5\). Indeed, if
    \(5\in I\), then
    \[
        5=(-\sqrt{-5})(\sqrt{-5})\in I
    \]
    and primality would force \(\sqrt{-5}\in I\), again making
    \(2b\sqrt{-5}\in I\). Hence \(\gcd(p,10b)=1\).

    Choose \(x,y\in\integers\) such that
    \[
        -10bx+py=1.
    \]
    Modulo \(I\),
    \begin{align*}
        (\alpha+I)(x\sqrt{-5}+I)
        &=
        (2b\sqrt{-5}+I)(x\sqrt{-5}+I)\\
        &=
        -10bx+I\\
        &=
        1+I.
    \end{align*}
    Thus every nonzero element of \(A/I\) is invertible in both cases.
    Therefore \(A/I\) is a \field and \(I\) is maximal.
\end{snippetproof}

\begin{snippet}{sqrt-minus-five-zero-prime-ideal-not-maximal}
    The zero \ideal of \(\integers[\sqrt{-5}]\) is prime because the
    \ring is an \integraldomain, but it is not maximal. The nonzero
    hypothesis in the preceding proposition is therefore essential.
\end{snippet}

\end{document}
